% Generated draft: requires mathematical and source review.
\documentclass[11pt]{article}
\usepackage[margin=25mm]{geometry}
\usepackage{fontspec}
\setmainfont{Latin Modern Roman}
\usepackage{amsmath,amssymb,mathtools,unicode-math}
\setmathfont{Latin Modern Math}
\usepackage{enumitem,needspace,xurl,hyperref}
\setlength{\parindent}{0pt}
\setlength{\parskip}{5pt}
\setlength{\emergencystretch}{3em}
\newcommand{\N}{\mathbb{N}}
\newcommand{\Z}{\mathbb{Z}}
\newcommand{\Q}{\mathbb{Q}}
\newcommand{\R}{\mathbb{R}}
\newcommand{\C}{\mathbb{C}}
\newcommand{\F}{\mathbb{F}}
\newcommand{\A}{\mathbb{A}}
\newcommand{\PP}{\mathbb{P}}
\newcommand{\E}{\mathbb{E}}
\newcommand{\eps}{\varepsilon}
\newcommand{\dd}{\,\mathrm{d}}
\begin{document}
\section*{UnsolvedMath ID: 30005622}
\textbf{Source problem code:} catalogue entry
\subsection{Definitions and mathematical statement}
Let \(K\in\R\setminus\{0\}\). For a Riemannian surface \((M,g)\) with boundary and sufficient regularity for a discrete Neumann spectrum, write
\[
0=\mu_0(M,g)\leq\mu_1(M,g)\leq\mu_2(M,g)\leq\cdots,
\]
counting multiplicities, for the eigenvalues of
\[
-\Delta_g u=\mu u\quad\text{in }M,\qquad
\partial_\nu u=0\quad\text{on }\partial M,
\]
where \(\nu\) is the outward unit normal. Under this indexing convention, the third eigenvalue is \(\mu_2\). Write \(\kappa_g\) for Gaussian curvature. Whenever such a disk exists, let \(D_K(a)\) denote a geodesic disk of area \(a\) in the complete simply connected surface of constant Gaussian curvature \(K\).

The supplied source asks whether, among simply connected surfaces satisfying \(\kappa_g\leq K\), the third Neumann eigenvalue is maximized by a disjoint union of two constant-curvature-\(K\) disks. Essential specifications are missing: the source does not state an area constraint or another normalization, the relative sizes of the disks, the boundary regularity, or the indexing convention. Moreover, the proposed disconnected maximizer is not itself a simply connected surface under the usual convention that simply connected spaces are connected.

Under the additional, unconfirmed interpretation of fixed total area \(A>0\), two equal-area comparison disks, and the indexing above, a candidate precise inequality is
\[
\mu_2(M,g)\leq
\mu_2\!\left(D_K(A/2)\sqcup D_K(A/2)\right)
\]
for every admissible simply connected \((M,g)\) satisfying
\[
\operatorname{Area}_g(M)=A,\qquad \kappa_g(x)\leq K\quad(x\in M).
\]
This conditional formulation applies only when the comparison disks exist; the intended admissible area range for \(K>0\) is not supplied. It is also unresolved whether the source intends attainment in an enlarged disconnected class or a sharp supremum over connected surfaces. These additional interpretations require confirmation before this can be treated as a fully specified conjecture.

The original passage states that the curvature-bound conjecture is known for \(K=0\) by Girouard and Polterovich [10], and open for \(K<0\) and \(K>0\). It separately attributes constant-curvature results to Bucur and Henrot [3] in the Euclidean setting, Freitas and Laugesen [9] in the hyperbolic setting, and Bucur, Martinet and Nahon [4] for spherical domains. These are supplied source claims, not independently verified mathematical or literature conclusions.
\subsection{Short English statement}
The question concerns whether two constant-curvature-\(K\) disks give the sharp upper bound for the third Neumann eigenvalue of simply connected surfaces with Gaussian curvature at most \(K\), for \(K\neq0\). A precise conjecture still requires a normalization, disk sizes, admissible surface and area conditions, and clarification of how a disconnected comparison realizes the maximum. The displayed fixed-area, equal-disk formulation is explicitly conditional. The supplied record labels the problem partially solved.
\subsection{Sources}
\begin{itemize}
\item Local UnsolvedMath registry, \texttt{data/unsolvedmath/problems.json}; permanent ID \(30005622\), source problem code \texttt{catalogue entry}, title ``Third Neumann Eigenvalue Maximization Under Curvature Bounds.'' The supplied record gives proposed year 2023 and no named proposer.
\item Supplied report citation: \emph{Geometric Spectral Theory (2024)}, Oberwolfach Reports, pp.\ 2017--2106, DOI \texttt{10.4171/owr/2023/36}. \url{https://doi.org/10.4171/owr/2023/36}. The record labels its statement \texttt{corrected\_verified} and says it was reconstructed and checked against the official report; that verification is supplied metadata, not verification performed here.
\item Original-passage attribution: Girouard and Polterovich [10] for the \(K=0\) curvature-bound case, and Bucur and Henrot [3] for the Euclidean constant-curvature case. The supplied record does not provide their full bibliographic entries or individual links.
\item Supplied literature source: Freitas--Laugesen, hyperbolic space, corresponding to Freitas and Laugesen [9] in the original passage. \url{https://doi.org/10.2422/2036-2145.202010_059}.
\item Supplied literature source: Bucur--Martinet--Nahon, sphere, corresponding to Bucur, Martinet and Nahon [4] in the original passage. \url{https://arxiv.org/abs/2208.11413}.
\item Literature assessment dated exactly \texttt{2026-08-22}: the supplied curation labels the problem partially solved and reports the two-disk maximizer as proved in Euclidean and hyperbolic constant-curvature settings and for spherical domains, while the broader upper-curvature-bound problem remains open. This is supplied dated curation derived from the OWR record, not independently verified here; its time-sensitive status and cited claims require independent checking.
\end{itemize}
\end{document}
